\documentclass[11pt,a4paper]{article}
\usepackage{gtu-el-yazisi}
\begin{document}
\notbilgi{Alan Teorisi}{Yüzeysel yüke etki eden kuvvet}{Ders11.pdf}

\sayfa{1}
\section{Yüzeysel yüke etki eden kuvvet}

$q_1 \ldots q_n$ tane yük olsun.

$q_1$ yüküne etki eden kuvveti:
\[ \vec F_1 = \frac{q_1}{4\pi\varepsilon}\sum_{i=2}^{n} q_i\,\frac{\vec r_i}{r_i^{3}} \]

\sekilnotu[1]{$V$ bölgesi (kapalı oval bir bölge) içinde $\rho_v$ yük yoğunluğuyla dağıtılmış yükler. Bölge içinde bir $P$ noktası ve etrafında ``$S_\mu$ küredir'' diye etiketlenmiş küçük bir küre (turuncu/kırmızı daire, yarıçap $\mu$); $P$'den bölgedeki bir yük elemanına $r'$ vektörü, yanında $\rho_v$. Sağda $V$'yi gösteren ok: ``yük yoğunluğu ile yükleri dağıtalım''; $P$'den aşağı ok: ``$P$ noktasına etki eden kuvveti hesaplayalım''.}

$P$ noktasına etki eden kuvveti hesaplayalım.

$S_\mu$ küresi dışında bulunan bütün yükleri içerideki birim yüke etki kuvvet (Elektrik alan). Bu kuvveti hesaplayıp, $\mu\to 0$ götürülürse $P$ noktasındaki elektrik alanı buluruz.
\[ \vec E(P) = \frac{1}{4\pi\varepsilon}\lim_{\mu\to 0}\int_{V-V_\mu} \rho_v\,\frac{\vec r'}{r'^{3}}\,dV \]

\sekilnotu[1]{Eğri bir $S$ yüzeyi (``$S$ yüzeyi'', ``$\rho_s$ yoğunluğu'') ve üzerinde $P$ noktası. $P$'nin hemen üstünde $P^+$ (turuncu $\vec n$ ``normal vektör'' yönünde, mesafe $h$), hemen altında $P^-$ (mesafe $h$). $P$ çevresinde mavi kesikli küçük disk, ``$S_\mu$ yüzeyi'', yarıçap $\mu$. Yüzeyin sol alt ucunda bir $ds$ elemanı; $ds$'den $P^+$, $P$ ve $P^-$ noktalarına turkuaz $\vec r_+$, $\vec r$, $\vec r_-$ vektörleri. Sağda ayrı çizim ``$S_\mu$ nedir'': yüzeyi dik kesen bir silindir (kutu), yarıçapı $\mu$, yüzeyle kesişen taralı disk $S_\mu$.}

$P$ noktası üzerindeki elektrik alanı bulalım.
\[ \vec E(P) = \frac{1}{4\pi\varepsilon}\lim_{\mu\to 0}\int_{S-S_\mu} \rho_s\,\frac{\vec r'}{r'^{3}}\,ds \]

Daha basit bir yol bulmaya çalışalım
\begin{align*}
\vec E(P_+) &= \frac{1}{4\pi\varepsilon}\int_S \rho_s\,\frac{\vec r_+}{r_+^{3}}\,ds \\
\vec E(P_-) &= \frac{1}{4\pi\varepsilon}\int_S \rho_s\,\frac{\vec r_-}{r_-^{3}}\,ds
\end{align*}

\sayfa{2}
\begin{align*}
\vec E(P_+) &= \frac{1}{4\pi\varepsilon}\int_{S-S_\mu}\rho_s\,\frac{\vec r_+}{r_+^{3}}\,ds
  + \overbrace{\frac{1}{4\pi\varepsilon}\int_{S_\mu}\rho_s\,\frac{\vec r_+}{r_+^{3}}\,ds}^{e(h,\mu)\,\hat n} \\[4mm]
\vec E(P_-) &= \frac{1}{4\pi\varepsilon}\int_{S-S_\mu}\rho_s\,\frac{\vec r_-}{r_-^{3}}\,ds
  + \underbrace{\frac{1}{4\pi\varepsilon}\int_{S_\mu}\rho_s\,\frac{\vec r_-}{r_-^{3}}\,ds}_{-e(h,\mu)\,\hat n}
\end{align*}

(yani $h\to 0$'a götürsek)
\begin{align*}
\lim_{P_+\to P}\vec E(P_+) &= \vec E(P) + e(0,\mu)\,\hat n \\
\lim_{P_-\to P}\vec E(P_-) &= \vec E(P) - e(0,\mu)\,\hat n
\end{align*}
\begin{align*}
\vec E^{+}(P) &= \vec E(P) + e(0,\mu)\,\hat n \\
\vec E^{-}(P) &= \vec E(P) - e(0,\mu)\,\hat n
\end{align*}

$\mu\to 0$'a götürürsek
\[
\begin{array}{rrcl}
 & \vec E^{+}(P) &=& \vec E(P) + e(0)\,\hat n \\
+\ & \vec E^{-}(P) &=& \vec E(P) - e(0)\,\hat n \\ \hline
 & \vec E(P) &=& \tfrac12\bigl[\vec E^{+}(P)+\vec E^{-}(P)\bigr]
\end{array}
\]

Herhangi bir $\rho_s\,ds$ yüzeysel yüküne etki eden kuvvet:
\[ d\vec F = \frac12\,\rho_s\,ds\,\bigl[\vec E^{+}+\vec E^{-}\bigr] \]

\subsection{Sayfa 55 problem 1}
La la la laaa la la laaaa

\begin{center}
\begin{tikzpicture}[line width=0.8pt,murekkep]
  \draw[-{Stealth}] (0,-0.7) -- (0,2.3) node[right] {$z$};
  \draw (-1.5,1.4) -- (5.2,1.4);
  \draw (-1.5,0) -- (5.2,0);
  \node[anchor=west] at (0.05,1.0) {$h$};
  \node[anchor=west] at (0.05,-0.35) {$0$};
  \node at (4.75,1.0) {$\rho_h$};
  \node at (4.9,-0.4) {$\rho_0$};
  \draw[-{Stealth}] (5.2,0) .. controls (5.65,0.5) and (5.85,1.3) .. (5.95,1.75) node[right] {yüzey};
\end{tikzpicture}
\end{center}

a) $z=0$'daki ve $z=h$'taki elektrik alanı doğrudan bulunuz.

\hspace{1.5em}{\small sadece $\rho_h$'ın etkisi var diyelim}
\[ \vec E(0) = -\frac{\rho_h}{2\varepsilon}\,\vec e_z \qquad\qquad \vec E(h) = +\frac{\rho_0}{2\varepsilon}\,\vec e_z \]

b) $z>h$, \ $0<z<h$, \ $z<0$
\[ \vec E = \frac{\rho_0+\rho_h}{2\varepsilon}\,\vec e_z \]

\sayfa{3}
$0<z<h$
\[ \vec E = \frac{(\rho_0-\rho_h)}{2\varepsilon}\,\vec e_z \]
$z<0$
\[ \vec E = -\frac{(\rho_0+\rho_h)}{2\varepsilon}\,\vec e_z \]

c) $z=0$ ve $z=h$

\noindent
\begin{minipage}[t]{0.49\linewidth}
\begin{align*}
\vec E(0) &= \tfrac12\bigl[\vec E^{+}(0)+\vec E^{-}(0)\bigr] \\
\vec E^{+}(0) &= \frac{(\rho_0-\rho_h)}{2\varepsilon}\,\vec e_z \\
\vec E^{-}(0) &= -\frac{(\rho_0+\rho_h)}{2\varepsilon}\,\vec e_z \\
\vec E(0) &= \frac12\left[\frac{\rho_0-\rho_h-\rho_0-\rho_h}{2\varepsilon}\,\vec e_z\right] \\
 &= -\frac{\rho_h}{2\varepsilon}\,\vec e_z
\end{align*}
\end{minipage}\hfill
\begin{minipage}[t]{0.47\linewidth}
\[
\left\{\begin{aligned}
\vec E(h) &= \tfrac12\bigl[\vec E^{+}(h)+\vec E^{-}(h)\bigr] \\
\vec E^{+}(h) &= \frac{\rho_0+\rho_h}{2\varepsilon}\,\vec e_z \\
\vec E^{-}(h) &= \frac{(\rho_0-\rho_h)}{2\varepsilon}\,\vec e_z \\
\vec E(h) &= \frac12\left[\frac{\rho_0+\rho_h+\rho_0-\rho_h}{2\varepsilon}\right]\vec e_z \\
\vec E(h) &= \frac{\rho_0}{2\varepsilon}\,\vec e_z
\end{aligned}\right.
\]
\end{minipage}

d) $z=0$ üzerinde $S$ alan parçasına etki eden kuvvet
\begin{align*}
\vec F &= \frac12\,\rho_0\,S\,\bigl[\vec E^{+}(0)+\vec E^{-}(0)\bigr] \\
 &= -\frac{\rho_0\,\rho_h\,S}{2\varepsilon}\,\vec e_z
\end{align*}

e) $\rho_h=\rho_0 \;\Rightarrow\; \vec F = -\dfrac{\rho_0^{2}}{2\varepsilon}\,\vec e_z$

\hspace{1.5em}$\rho_h=-\rho_0 \;\Rightarrow\; \vec F = \dfrac{\rho_0^{2}}{2\varepsilon}\,\vec e_z$

\end{document}
